\section{Behavioral coverage and peer expressibility}
\label{sec:framework}

Let $\mathcal J=\{1,\ldots,N\}$ index a model collection, $t$ a target, and $\mathcal J_{-t}$ its peers. An audit observes the models on the same transformed input $T_{\theta,\omega}(x)$, where $\theta$ specifies an intervention level and $\omega$ its realization. Writing $z=(x,\theta,\omega)$, the recorded response is
\begin{equation}
Y_j(z)=g_x\!\left(f_j(T_{\theta,\omega}(x))\right)\in\R^q.
\label{eq:response}
\end{equation}
The response map $g_x$ selects what the audit compares. For multiple-choice questions, it can record reference-answer confidence or the entire option-probability vector. For forecasting, it records the prediction. Model-specific wrappers preserve the same semantic input and recorded quantity across peers.

\paragraph{Fixed collective reconstruction.}
Let $\Delta^{N-2}=\{w\in\R^{N-1}:w_j\geq0,\sum_jw_j=1\}$. The peer-response class is
\begin{equation}
h_w(z)=\sum_{j\in\mathcal J_{-t}}w_jY_j(z),\qquad
\mathcal C_t=\{h_w:w\in\Delta^{N-2}\}.
\label{eq:peer-class}
\end{equation}
A single $w$ applies across the audit. Even when $Y_j(z)$ is scalar, the fit reconstructs a whole response trace $(Y_j(z_1),\ldots,Y_j(z_m))$, rather than choosing coefficients separately for each observation.

Convex weights combine the existing responses on their original scale and preserve probability normalization for vector responses. Allowing signed coefficients would permit extrapolation; allowing input-dependent weights would permit adaptive selection. The fixed convex class isolates the coverage supplied by an unchanged combination of peers.

A fitting distribution $P$ defines the population approximation,
\begin{equation}
h_t^\star\in\arg\min_{h\in\mathcal C_t}\E_P\|Y_t-h\|_2^2.
\label{eq:projection}
\end{equation}
For evaluation distribution $Q$, the residual and its mean magnitude are
\begin{equation}
R_t(z)=Y_t(z)-h_t^\star(z),\qquad
U_t(P,Q)=\E_Q\|R_t(Z)\|_2.
\label{eq:pier}
\end{equation}
We call $R_t$ \PIER and $U_t$ the \emph{coverage gap}. Scalar responses use absolute residuals; the full-option analysis additionally reports total variation. Section~\ref{sec:measurement} specifies the support conditions under which the fitted projection defines a unique evaluation response.

\begin{figure}[!t]
\centering
\begin{subfigure}[t]{0.485\linewidth}
\centering\includegraphics[width=\linewidth]{figures/pdf/fig1a_matched_audit.pdf}
\caption{Matched responses and a fixed peer fit.}
\end{subfigure}\hfill
\begin{subfigure}[t]{0.485\linewidth}
\centering\includegraphics[width=\linewidth]{figures/pdf/fig1b_peer_projection.pdf}
\caption{Projection onto the peer-response class.}
\end{subfigure}
\caption{\textbf{A fixed peer fit reconstructs a response trace.} A fixed combination is learned from matched fitting observations and evaluated on held-out inputs. Its residual records what remains unexplained by the peer class; comparing alternative peer sets identifies the source of coverage. Both panels are conceptual.}
\label{fig:framework}
\end{figure}

\paragraph{Estimation and peer dependence.}
\DISCO solves the empirical simplex-constrained version of Eq.~\eqref{eq:projection} on a fitting sample and evaluates $\widehat U_t$ on held-out observations. The audit requires model responses but no model parameters or gradients. Absolute-loss fitting supplies a second matched comparison in the scalar experiments; Section~\ref{sec:measurement}'s projection uniqueness result applies to squared loss.

Let $E_t^{\mathrm{single}}$ be the held-out error of the peer selected on the fitting sample, and $E_t^{\mathrm{group}}$ the error of the fitted combination. Define
\begin{align}
\Gamma_t&=E_t^{\mathrm{single}}-E_t^{\mathrm{group}},\label{eq:groupgain}\\
I_{t\leftarrow j}&=\widehat U_t(\mathcal J_{-t}\setminus\{j\})-\widehat U_t(\mathcal J_{-t}).\label{eq:influence}
\end{align}
Each reduced peer set is refitted under the same design. Positive $\Gamma_t$ identifies a collective advantage; positive $I_{t\leftarrow j}$ measures how much coverage is lost without peer $j$. Together they separate dependence on one peer from improvement supplied by the rest. They are held-out comparisons and can have either sign.
